\section{Enumerations, and what is machine-checked}\label{app:counts}

Every count the paper relies on is collected here with its provenance.
Three provenances occur, and the distinction matters because only the
first is a theorem: \emph{kernel} for a number a Lean declaration
computes or a Lean proof establishes; \emph{exhaustive} for a number
produced by a search outside the proof assistant, reproducible but not
verified; \emph{closed form} for a number read off an identity proved in
the body.  No argument in the body rests on a count of the second kind.

\subsection{Active-family census at
\texorpdfstring{$(6,3)$ and $(7,3)$}{(6,3) and (7,3)}}\label{app:census}

By \eqref{eq:coverage} the family $A^{+}$ of active $k$-subsets
carrying a positive multiplier covers the index set.  The covering families
of $3$-subsets are counted here up to the symmetric group of the index set
and listed by cardinality.  At $m=6$ a covering family has at
least two members and at most twenty; Table~\ref{tab:census} lists the
counts, whose total is $2102$.  The single class of cardinality two is the
partition branch that Corollary~\ref{cor:e3} leaves undecided; setting
it aside leaves $2101$.  By Proposition~\ref{prop:carath} the support of the
multiplier vector may be taken of size at most $15$ here, so only the columns
up to $15$ bear on
Conjecture~\ref{conj:sixthree}; they account for $2069$ of the $2102$ classes,
and for $2068$ of the $2101$ with at least three members.

At $m=7$ a covering family has at least three members, and the counts are
\[
\begin{array}{c|cccccccc}
\lvert A^{+}\rvert & 3&4&5&6&7&8&9&10\\ \hline
\text{classes} & 3&17&94&415&1599&5308&15397&39081\\[4pt]
\lvert A^{+}\rvert & 11&12&13&14&15&16&17&18\\ \hline
\text{classes} & 87347&172684&303116&473733&660907&824389&920439&920443\\[4pt]
\lvert A^{+}\rvert & 19&20&21&22&23&24&25&26\\ \hline
\text{classes} & 824409&660949&473827&303277&172933&87659&39433&15709\\[4pt]
\lvert A^{+}\rvert & 27&28&29&30&31&32&33&34\\ \hline
\text{classes} & 5557&1760&509&137&38&10&3&1
\end{array}
\]
together with a single class of cardinality $35$,
with total $7\,011\,184$.  By Proposition~\ref{prop:carath} the bound at
$(7,3)$ is $19$, and the classes of cardinality at most
$19$ already number $5\,249\,381$.  Enumeration at this size is not a route;
Problem~\ref{prob:seventhree} asks instead for an argument uniform in
the active family.
